\begin{prob}
    For each of the following functions, express its rate of growth using
    $\Theta$-notation, and prove your answer by finding constants which satisfy
    the definition of $\Theta$-notation. Make sure to show your work by writing
    out the chain of inequalities to prove each bound, like we did in lecture.
    Note: please do \emph{not} use limits to prove your answer (though you can
    use them to check your work).

    \begin{subprobset}
        \begin{subprob}
            $f(n) = 3n^4 + 10n^2 + 4$

            \begin{soln}
                $f(n) = \Theta(n^4)$.

                We start with an upper bound:
                \begin{align*}
                    f(n)
                    &=
                    3n^4 + 10n^2 + 4
                        \\
                    &\leq
                        3n^4 + 10n^4 + 4n^4 \qquad (n \geq 1)
                        \\
                    &=
                        17 n^4
                \end{align*}
                A simple lower bound is $f(n) \geq 3n^4$. Hence
                $3n^4 \leq 3n^4 + 10n^2 + 4 \leq 17n^4$ for all $n \geq 1$.
            \end{soln}

        \end{subprob}

        \begin{subprob}
            $
            \displaystyle
            f(n) =
            \frac{n^2 + 2n - 5}{n - 10}
            $

            \begin{soln}
                $f(n) = \Theta(n)$.

                We begin with an upper bound.
                \begin{align*}
                    f(n)
                    &=
                        \frac{n^2 + 2n - 5}{n - 10}
                        \\
                    &\leq
                        \frac{n^2 + 2n}{n - 10}
                        \\
                    &=
                        \frac{n(n + 2)}{n - 10}
                        \\
                        \intertext{Now, $(n+2)/(n-10) \leq 2$ for all $n \geq 22$, hence:}
                    &\leq
                        2 n, \qquad (n \geq 22)
                        \\
                \end{align*}
                Now for a lower bound:
                \begin{align*}
                    f(n)
                    &=
                        \frac{n^2 + 2n - 5}{n - 10}
                        \\
                    &\geq
                        \frac{n^2 - 5}{n - 10}
                        \\
                    \intertext{To get a lower bound, we make this
                    quantity smaller. We can do so by making the
                    denominator bigger by dropping the $-10$:}
                    &\geq
                        \frac{n^2 - 5}{n}
                        \\
                    &\geq
                    \frac{\frac{1}{2}n^2  + \frac{1}{2}n^2 - 5}{n}
                        \\
                        \intertext{Now, $\frac{1}{2}n^2 -5 \geq 0$ for all $n \geq 4$. Hence:}
                    &\geq
                \frac{\frac{1}{2}n^2  + 0}{n}, \qquad (n \geq 4)
                        \\
                    &= n/2
                \end{align*}
                So in total, $\frac{n}{2} \leq f(n) \leq 2n$ for all $n \geq 22$.
            \end{soln}
            
        \end{subprob}

        \begin{subprob}
            $
            \displaystyle
            f(n) = 
            \begin{cases}
                n + 10, & n \text{ is odd},\\
                n - 5, & n \text{ is even}
            \end{cases}
            $

            \begin{soln}
                $f(n) = \Theta(n)$.

                In either case, $f(n) \leq 2n$ for all $n \geq 10$.
                Likewise, $f(n) \geq n/2$ for all $n \geq 10$. Hence:
                $\frac{n}{2} \leq f(n) \leq 2n$ for all $n \geq 10$.
            \end{soln}
        \end{subprob}
    \end{subprobset}
\end{prob}
